\documentclass[11pt]{article}
\usepackage{preamble}
\setlength{\parskip}{4pt}

\begin{document}

\daybanner{AP Statistics \textperiodcentered\ Topic 1.2 (CED 1.2) \textperiodcentered\ B: 2026-09-10 \textperiodcentered\ E: 2026-09-11 \textperiodcentered\ TEACHER KEY}{Numbers That Aren't Quantities}

\sectionbanner{CED TETHER}

{\small
\begin{itemize}[leftmargin=1.5em,itemsep=2pt,topsep=2pt]
  \item \textbf{SKILL: 2.A} --- Describe data presented numerically or graphically.
  \item ENDURING UNDERSTANDING: VAR-1 Given that variation may be random or not, conclusions are uncertain.
  \item VAR-1.B Identify variables in a set of data. [Skill 2.A]
  \item VAR-1.B.1 A variable is a characteristic that changes from one individual to another.
  \item VAR-1.C Classify types of variables. [Skill 2.A]
  \item VAR-1.C.1 A categorical variable takes on values that are category names or group labels.
  \item VAR-1.C.2 A quantitative variable is one that takes on numerical values for a measured or counted quantity.
\end{itemize}
}
\textbf{Essential question:} When is a number a quantity, and when is it just a label?\par
\textbf{DOK-3 skill:} 4.B \textperiodcentered\ \textbf{FRQ pattern:} \texttt{number-coded-categorical-decide-and-name-the-meaningless-summary}\par
\textbf{DOK rationale:} Part (c) asks the student to adjudicate two students' classifications of number-coded variables, justify each with the definition, and predict what specific calculation becomes meaningless under the wrong choice; both classifications of the rating can be defended, so the reasoning, not the label, is what is scored.\par

\frameworkphaseheader{Do Now (first take) $\rightarrow$ Explore (video + follow-along) $\rightarrow$ Exit (finish + turn in)}{1 $\rightarrow$ 3}{45}{%
\begin{itemize}[leftmargin=1.5em,itemsep=2pt,topsep=2pt]
  \item Board slide up at the bell. Read Maya's and Jordan's lines aloud; do not react.
  \item Video follow-along from minute 5.
  \item Board slide back up at minute 33. Circulate; push for the definition, not the label.
  \item Collect at the door. Score part (c) only, E/P/I. Either position on Rating can earn E.
\end{itemize}
}{%
\begin{itemize}[leftmargin=1.5em,itemsep=2pt,topsep=2pt]
  \item One-sentence first take on the zip-code question before any teaching.
  \item Video follow-along (individuals, variables, categorical vs quantitative).
  \item Finish (a)--(c); exit line; turn in.
\end{itemize}
}{%
\begin{itemize}[leftmargin=1.5em,itemsep=2pt,topsep=2pt]
  \item What does a zip code measure? What does it count?
  \item If I averaged two zip codes, what place would I get?
  \item Is a 5-star rating twice as good as a 2.5-star rating? Does that matter for your answer?
  \item Which column would you sort to find the `biggest'? Which column has no `biggest'?
\end{itemize}
}{Solo teacher. Circulate ~5 pairs. Never say `categorical' or `quantitative' yourself during the finish window; ask what the number measures.}

\begin{callout}{\IconWarn\ WATCH FOR}{calloutred}
\begin{itemize}[leftmargin=1.5em,itemsep=2pt,topsep=2pt]
  \item `It's a number so it's quantitative' --- the digit-vs-quantity confusion the whole lesson is about.
  \item Students who get Zip code right by instinct but cannot state the definition (P, not E).
  \item Listing the newspaper or the school as the individual.
\end{itemize}
\end{callout}

\sectionbanner{SCORING PART (c) --- E / P / I}

\textbf{Expected elements}
\begin{itemize}[leftmargin=1.5em,itemsep=2pt,topsep=2pt]
  \item \textbf{zip-categorical-with-definition} (required) --- Says zip code is categorical AND ties it to the definition (measured or counted quantity vs. a label)
  \item \textbf{rating-decided-with-reason} (required) --- Takes a position on Rating and defends it with the label-vs-quantity distinction (either position accepted with a correct reason)
  \item \textbf{names-meaningless-calc} (required) --- Names a specific calculation (mean, SD, histogram of zip codes) and explains why it has no meaning
\end{itemize}

{\small\renewcommand{\arraystretch}{1.25}
\noindent\begin{tabular}{|p{0.5in}|p{5.9in}|}
\hline
\textbf{E} & Zip code categorical with the definition stated, a defended position on Rating, and a specific meaningless calculation explained (e.g. mean zip code is not a place). \\ \hline
\textbf{P} & Two of the three elements, OR all three present but the zip-code reason is only `it's not really a number' with no definition, OR the meaningless calculation is named but not explained. \\ \hline
\textbf{I} & Agrees with Maya on zip code, or gives no definition-based reasoning, or names no calculation. \\ \hline
\end{tabular}}

\textbf{Common mistakes}
\begin{itemize}[leftmargin=1.5em,itemsep=2pt,topsep=2pt]
  \item `It has digits, so it's quantitative' --- confuses how a value is written with what it measures
  \item Saying zip code is categorical `because you can't add them' without saying WHY (it measures/counts nothing)
  \item Treating Rating as obviously quantitative because the quiz app computes an average star rating
  \item Naming the individuals as `the newspaper' or `the school' instead of the six students
\end{itemize}

\bankitem{Optional \#1 --- not collected}{\dokbadge{2}~For each variable, write \textbf{C} (categorical) or \textbf{Q} (quantitative); for each Q, give the unit: (1) jersey number of a soccer player; (2) number of goals scored this season; (3) shoe size; (4) preferred position (goalkeeper, defender, midfielder, forward); (5) minutes played in the last match.}{}
\answer{(1) C --- a label, even though numeric. (2) Q, goals. (3) Q, shoe-size units (a measured foot length on a scale); accept C with a reason that sizes are ordered labels. (4) C. (5) Q, minutes.}

\clearpage
\focusbanner{TODAY'S PROBLEM --- annotated}

A school newspaper surveyed six students for a story about getting to school. Part of the data table is shown.\par\smallskip Two reporters disagree about the numbers in it. \textbf{Maya:} ``Zip code and Rating are both quantitative --- they're numbers, so we can average them.'' \textbf{Jordan:} ``They're both categorical. Numbers can be labels.''\par
\begin{center}
\textbf{Getting-to-school survey (six students)}\par\smallskip
\begin{tabular}{|l|c|c|c|c|}
\hline
\textbf{Student} & \textbf{Lunch period} & \textbf{Zip code} & \textbf{Commute (min)} & \textbf{Rating (1--5 stars)} \\ \hline
\textbf{Aiden} & A & 01902 & 12 & 4 \\ \hline
\textbf{Bea} & B & 01905 & 35 & 2 \\ \hline
\textbf{Cruz} & A & 01902 & 8 & 5 \\ \hline
\textbf{Dev} & C & 01904 & 22 & 3 \\ \hline
\textbf{Elise} & B & 01905 & 41 & 1 \\ \hline
\textbf{Fiona} & A & 01901 & 15 & 4 \\ \hline
\end{tabular}
\end{center}
\begin{firsttakebox}{FIRST TAKE (5 min)}
Is a zip code a quantitative variable? One sentence: yes or no, and the reason you'd give Maya.\par
\answer{Not graded. Most will say `yes, it's a number'; that is the hook.}
\end{firsttakebox}

\textit{Rules callout ``VARIABLES (keep on the page)'' is printed on the student sheet.}\par\medskip

\dokbadge{1}~\textbf{(a)} Who are the individuals in this data set, and what are the four variables?\par
\answer{Individuals: the six students (Aiden, Bea, Cruz, Dev, Elise, Fiona). Variables: lunch period, zip code, commute time, rating of the morning commute.}

\dokbadge{2}~\textbf{(b)} Classify each of the four variables as categorical or quantitative. For each quantitative variable, state its unit.\par
\answer{Lunch period: categorical. Zip code: categorical (numeric labels for places; arithmetic on them is meaningless). Commute time: quantitative, minutes. Rating (1--5 stars): categorical --- the numbers are ordered labels for opinions; accept quantitative ONLY if the student states that the numbers are being treated as a measured score and names the unit as `stars'.}

\dokbadge{3}~\textbf{(c)} Decide who is right about \textbf{Zip code} and who is right about \textbf{Rating}, using the definition of a quantitative variable in each case. Then name one specific calculation a reader would find meaningless if Maya's classification were used for Zip code, and explain why it is meaningless.\par
\answer{Zip code: Jordan is right. A quantitative variable takes numerical values for a measured or counted quantity (VAR-1.C.2); a zip code measures nothing and counts nothing --- it names a place, so it is categorical even though it is written with digits. Rating: Jordan is right under the CED definition (the values 1--5 are category labels for opinions, VAR-1.C.1), though a student who argues that the rating is being used as a measured score and says so explicitly can defend `quantitative'. Meaningless calculation for zip code under Maya's view: the mean zip code (e.g. averaging 01902 and 01905 to get 01903.5) --- it is not a place, and `larger zip code' does not mean `more of anything'. Also acceptable: a histogram or standard deviation of zip codes.}

\begin{summaryexitbox}{EXIT}
One thing the video changed about my first take: \hrulefill
\end{summaryexitbox}

\end{document}
